% =====================================================================
%  PREÁMBULO DE LA TESIS: paquetes y formato
%  Ajuste aquí márgenes, letra e interlineado si el Reglamento de
%  Graduación de la Carrera de Ingeniería Civil (UATF) exige otros valores.
% =====================================================================

% --- Idioma, codificación y letra ------------------------------------
\usepackage[T1]{fontenc}
\usepackage[spanish,es-tabla,es-noshorthands,es-lcroman]{babel}
\usepackage{amsmath}
\usepackage{newtxtext,newtxmath}      % letra tipo Times New Roman
\renewcommand{\ttdefault}{lmtt}       % monoespaciada (rutas y DOI): Latin Modern
% El mapa de fuentes de MiKTeX de este equipo no registra newtx (MiKTeX
% indica "User/administrator updates are out-of-sync"); se carga aquí:
\pdfmapfile{+newtx.map}
\pdfmapfile{+qtm.map}   % TeX Gyre Termes (símbolos TS1 de newtxtext)
\pdfmapfile{+qhv.map}   % TeX Gyre Heros (sans)
\usepackage[protrusion=true,expansion=false]{microtype}
\usepackage[autostyle]{csquotes}

% --- Página: tamaño carta, márgenes e interlineado --------------------
\usepackage[letterpaper,top=2.5cm,bottom=2.5cm,left=3cm,right=2.5cm,
            headheight=15pt,footskip=1.2cm]{geometry}
\usepackage{setspace}
\onehalfspacing                         % interlineado 1,5
\setlength{\parindent}{0pt}
\setlength{\parskip}{6pt plus 2pt minus 1pt}

% --- Colores, cajas y utilidades --------------------------------------
\usepackage[dvipsnames,table]{xcolor}
\usepackage{etoolbox}
\usepackage[inline]{enumitem}
\setlist{itemsep=2pt,topsep=4pt,parsep=0pt,leftmargin=*}

% --- Títulos de capítulos y secciones ---------------------------------
\usepackage{titlesec}
\titleformat{\chapter}[display]
  {\normalfont\large\bfseries\centering}
  {\MakeUppercase{\chaptertitlename}\ \thechapter}{6pt}{\MakeUppercase}
\titleformat{name=\chapter,numberless}[block]
  {\normalfont\large\bfseries\centering}{}{0pt}{\MakeUppercase}
\titlespacing*{\chapter}{0pt}{0pt}{18pt}
\titleformat{\section}{\normalfont\normalsize\bfseries}{\thesection.}{0.6em}{}
\titleformat{\subsection}{\normalfont\normalsize\bfseries}{\thesubsection.}{0.6em}{}
\titleformat{\subsubsection}{\normalfont\normalsize\bfseries\itshape}{\thesubsubsection.}{0.6em}{}
\titlespacing*{\section}{0pt}{12pt}{4pt}
\titlespacing*{\subsection}{0pt}{10pt}{3pt}
\titlespacing*{\subsubsection}{0pt}{8pt}{2pt}
\setcounter{secnumdepth}{3}
\setcounter{tocdepth}{2}

% --- Nombres en español ------------------------------------------------
\addto\captionsspanish{%
  \renewcommand{\contentsname}{Índice general}%
  \renewcommand{\listtablename}{Índice de tablas}%
  \renewcommand{\listfigurename}{Índice de figuras}%
  \renewcommand{\appendixname}{Anexo}%
  \renewcommand{\bibname}{Bibliografía}%
}

% --- Índices (tabla de contenido) -------------------------------------
\usepackage{tocloft}
\renewcommand{\cftchapfont}{\bfseries}
\renewcommand{\cftchappagefont}{\bfseries}
\renewcommand{\cftchapleader}{\cftdotfill{\cftdotsep}}
\renewcommand{\cftchappresnum}{\chaptername\ }
\renewcommand{\cftchapaftersnum}{.}
\setlength{\cftchapnumwidth}{6.4em}
\setlength{\cftbeforechapskip}{6pt}
\renewcommand{\cfttabpresnum}{Tabla\ }
\renewcommand{\cftfigpresnum}{Figura\ }
\renewcommand{\cfttabaftersnum}{.}
\renewcommand{\cftfigaftersnum}{.}
\setlength{\cfttabnumwidth}{5.2em}
\setlength{\cftfignumwidth}{5.6em}
\setlength{\cfttabindent}{0pt}
\setlength{\cftfigindent}{0pt}

% --- Tablas ------------------------------------------------------------
\usepackage{array,booktabs,tabularx,longtable,multirow,makecell}
\renewcommand{\arraystretch}{1.15}
\newcolumntype{L}{>{\raggedright\arraybackslash}X}
\newcolumntype{C}{>{\centering\arraybackslash}X}
\newcolumntype{P}[1]{>{\raggedright\arraybackslash}p{#1}}
\newcolumntype{Q}[1]{>{\centering\arraybackslash}p{#1}}
\AtBeginEnvironment{table}{\small\setstretch{1}}
\AtBeginEnvironment{figure}{\setstretch{1}}

% --- Leyendas de tablas y figuras (estilo APA 7) -----------------------
%     Número en negrita y título en cursiva en la línea siguiente.
\usepackage{caption}
\captionsetup{font={small,stretch=1},labelfont=bf,textfont=it,
              labelsep=newline,justification=raggedright,
              singlelinecheck=false,skip=4pt}
\counterwithout{table}{chapter}         % Tabla 1, 2, 3...
\counterwithout{figure}{chapter}        % Figura 1, 2, 3...
\counterwithout{equation}{chapter}      % (1), (2), (3)...
% Nota debajo de una tabla o figura: \nota{Elaborado con base en ...}
\newcommand{\nota}[1]{\par\vspace{2pt}{\small\setstretch{1}\textit{Nota.} #1\par}}

% --- Figuras, planos y diagramas -----------------------------------------
\usepackage{graphicx}
\graphicspath{{figuras/}}
\usepackage{float}
\usepackage{pdflscape}
\usepackage{pdfpages}
\usepackage{tikz}
\usetikzlibrary{arrows.meta,positioning,babel}

% --- Unidades y números (coma decimal, punto de miles) ------------------
\usepackage{siunitx}
\sisetup{output-decimal-marker={,},group-separator={.},
         group-minimum-digits=4,per-mode=symbol}
\newcommand{\mdos}{m\textsuperscript{2}}    % metro cuadrado
\newcommand{\mtres}{m\textsuperscript{3}}   % metro cúbico

% --- Encabezado y pie de página --------------------------------------
\usepackage{fancyhdr}
\fancypagestyle{tesis}{\fancyhf{}\fancyfoot[C]{\thepage}%
  \renewcommand{\headrulewidth}{0pt}}
\fancypagestyle{plain}{\fancyhf{}\fancyfoot[C]{\thepage}%
  \renewcommand{\headrulewidth}{0pt}}
\pagestyle{tesis}

% --- Marcas de borrador -------------------------------------------------
%  \pendiente{...} : texto o dato que falta completar (rojo).
%  \guia{...}      : instrucciones para redactar una sección (recuadro gris).
%  Ambas desaparecen al poner \borradorfalse en tesis.tex.
\newif\ifborrador
\newcommand{\pendiente}[1]{\ifborrador{\color{BrickRed}\textbf{[Pendiente:} #1\textbf{]}}\fi}
\newcommand{\guia}[1]{\ifborrador\par\noindent
  \fcolorbox{gray!60}{gray!8}{\parbox{\dimexpr\linewidth-2\fboxsep-2\fboxrule\relax}{%
  \small\setstretch{1}\textbf{Guía de desarrollo.} #1}}\par\fi}

% --- Ruta a los datos de obra del repositorio -----------------------------
\newcommand{\datos}{../03_Datos_Obra_Edificio_Posgrado}

% --- Bibliografía APA 7.ª edición -----------------------------------------
\usepackage[style=apa,backend=biber]{biblatex}
\addbibresource{bibliografia.bib}

% --- Hipervínculos (cargar al final) --------------------------------------
\usepackage[hidelinks,bookmarksnumbered]{hyperref}
